
    %%%%%%%%%%%%%%%%%%%%%%%%%%%%%%%%%%%%%%%%%%%%%%%%%%%%%%%%%%%%%%%%%%%%%%%%%%%%%%%%
    %%% Classe do documento
    \documentclass[12pt, addpoints]{exam}
    \UseRawInputEncoding

    %%%%%%%%%%%%%%%%%%%%%%%%%%%%%%%%%%%%%%%%%%%%%%%%%%%%%%%%%%%%%%%%%%%%%%%%%%%%%%%%
    % preambulo.tex
    %
    % Arquivo de configuração de packages para uso o LaTeX, conforme minhas
    % preferências e modelos pessoais.
    %
    % Para maiores informações, visite:
    %    https://github.com/abrantesasf/latex
    %
    % NÃO ALTERE SE NÃO SOUBER O QUE ESTÁ FAZENDO!
    %%%%%%%%%%%%%%%%%%%%%%%%%%%%%%%%%%%%%%%%%%%%%%%%%%%%%%%%%%%%%%%%%%%%%%%%%%%%%%%%


    %%%%%%%%%%%%%%%%%%%%%%%%%%%%%%%%%%%%%%%%%%%%%%%%%%%%%%%%%%%%%%%%%%%%%%%%%%%%%%%%
    %%% Estruturas de controle:
    \input{/app/static/utils/controles}

    %%%%%%%%%%%%%%%%%%%%%%%%%%%%%%%%%%%%%%%%%%%%%%%%%%%%%%%%%%%%%%%%%%%%%%%%%%%%%%%%
    %%% Compilação condicional em PDF:
    \input{/app/static/utils/pdf}

    %%%%%%%%%%%%%%%%%%%%%%%%%%%%%%%%%%%%%%%%%%%%%%%%%%%%%%%%%%%%%%%%%%%%%%%%%%%%%%%%
    %%% Configurações de layout da página:
    \input{/app/static/utils/layout}

    %%%%%%%%%%%%%%%%%%%%%%%%%%%%%%%%%%%%%%%%%%%%%%%%%%%%%%%%%%%%%%%%%%%%%%%%%%%%%%%%
    %%% Configurações para autores:
    \input{/app/static/utils/autores}

    %%%%%%%%%%%%%%%%%%%%%%%%%%%%%%%%%%%%%%%%%%%%%%%%%%%%%%%%%%%%%%%%%%%%%%%%%%%%%%%%
    %%% Cabeçalho e rodapé:
    %%% Atenção! É MUITO DIFÍCIL ajustar apropriadamente os running headers
    %%% e running footers da primeira vez. Você pode confiar no LaTeX e deixar que
    %%% ele faça o ajuste automaticamente ou pode quebrar a cabeça e
    %%% tentar melhorar algo que o LaTeX já faz muito bem, mesmo que não fique
    %%% exatamente do jeito que você quer. O que você prefere? Se quiser ajustar
    %%% manualmente, descomente a linha abaixo e faça as configurações no arquivo
    %%% "utils/headings.tex".
    %\input{/app/static/utils/headings}

    %%%%%%%%%%%%%%%%%%%%%%%%%%%%%%%%%%%%%%%%%%%%%%%%%%%%%%%%%%%%%%%%%%%%%%%%%%%%%%%%
    %%% Configuração para as fontes do documento:
    %%% Se você quiser usar fontes próprias ou do sistema, precisará ajustar as
    %%% configurações no arquivo "utils/fontes.tex".
    %%% ATENÇÃO: por padrão eu utilizo XeLaTeX com fontes proprietárias projetadas
    %%% por Matthew Butterick, adquiridas em https://mbtype.com, que não podem ser
    %%% distribuídas devido à licença de utilização. Se você usar XeLaTeX, você
    %%% DEVERÁ OBRIGATORIAMENTE ajustar as fontes no arquivo "utils/fontes.tex".
    \input{/app/static/utils/fontes}

    %%%%%%%%%%%%%%%%%%%%%%%%%%%%%%%%%%%%%%%%%%%%%%%%%%%%%%%%%%%%%%%%%%%%%%%%%%%%%%%%
    %%% Sumário:
    \input{/app/static/utils/toc}

    %%%%%%%%%%%%%%%%%%%%%%%%%%%%%%%%%%%%%%%%%%%%%%%%%%%%%%%%%%%%%%%%%%%%%%%%%%%%%%%%
    %%% Matemática:
    \input{/app/static/utils/matematica}

    %%%%%%%%%%%%%%%%%%%%%%%%%%%%%%%%%%%%%%%%%%%%%%%%%%%%%%%%%%%%%%%%%%%%%%%%%%%%%%%%
    %%% Notas de rodapé:
    \input{/app/static/utils/footnotes}

    %%%%%%%%%%%%%%%%%%%%%%%%%%%%%%%%%%%%%%%%%%%%%%%%%%%%%%%%%%%%%%%%%%%%%%%%%%%%%%%%
    %%% Referências cruzadas, links, citações:
    \input{/app/static/utils/crossrefs}

    %%%%%%%%%%%%%%%%%%%%%%%%%%%%%%%%%%%%%%%%%%%%%%%%%%%%%%%%%%%%%%%%%%%%%%%%%%%%%%%%
    %%% Configurações para os hiperlinks em PDF:
    \input{/app/static/utils/pdfconfig}

    %%%%%%%%%%%%%%%%%%%%%%%%%%%%%%%%%%%%%%%%%%%%%%%%%%%%%%%%%%%%%%%%%%%%%%%%%%%%%%%%
    %%% Glossário e índice remissivo:
    \input{/app/static/utils/glossindex}

    %%%%%%%%%%%%%%%%%%%%%%%%%%%%%%%%%%%%%%%%%%%%%%%%%%%%%%%%%%%%%%%%%%%%%%%%%%%%%%%%
    %%% Referências bibliográficas:
    \input{/app/static/utils/bibliografia}

    %%%%%%%%%%%%%%%%%%%%%%%%%%%%%%%%%%%%%%%%%%%%%%%%%%%%%%%%%%%%%%%%%%%%%%%%%%%%%%%%
    %%% Cores:
    \input{/app/static/utils/cores}

    %%%%%%%%%%%%%%%%%%%%%%%%%%%%%%%%%%%%%%%%%%%%%%%%%%%%%%%%%%%%%%%%%%%%%%%%%%%%%%%%
    %%% Computação:
    \input{/app/static/utils/computacao}

    %%%%%%%%%%%%%%%%%%%%%%%%%%%%%%%%%%%%%%%%%%%%%%%%%%%%%%%%%%%%%%%%%%%%%%%%%%%%%%%%
    %%% Gráficos:
    \input{/app/static/utils/graficos}

    %%%%%%%%%%%%%%%%%%%%%%%%%%%%%%%%%%%%%%%%%%%%%%%%%%%%%%%%%%%%%%%%%%%%%%%%%%%%%%%%
    %%% Ícones:
    \input{/app/static/utils/icones}

    %%%%%%%%%%%%%%%%%%%%%%%%%%%%%%%%%%%%%%%%%%%%%%%%%%%%%%%%%%%%%%%%%%%%%%%%%%%%%%%%
    %%% Blocos:
    \input{/app/static/utils/blocos}

    %%%%%%%%%%%%%%%%%%%%%%%%%%%%%%%%%%%%%%%%%%%%%%%%%%%%%%%%%%%%%%%%%%%%%%%%%%%%%%%%
    %%% Boxes coloridos:
    \input{/app/static/utils/colorbox}

    %%%%%%%%%%%%%%%%%%%%%%%%%%%%%%%%%%%%%%%%%%%%%%%%%%%%%%%%%%%%%%%%%%%%%%%%%%%%%%%%
    %%% Tabelas:
    \input{/app/static/utils/tabelas}

    %%%%%%%%%%%%%%%%%%%%%%%%%%%%%%%%%%%%%%%%%%%%%%%%%%%%%%%%%%%%%%%%%%%%%%%%%%%%%%%%
    %%% Ambientes de listas:
    \input{/app/static/utils/listas}

    %%%%%%%%%%%%%%%%%%%%%%%%%%%%%%%%%%%%%%%%%%%%%%%%%%%%%%%%%%%%%%%%%%%%%%%%%%%%%%%%
    %%% Ambientes floats e similares:
    \input{/app/static/utils/floats}

    %%%%%%%%%%%%%%%%%%%%%%%%%%%%%%%%%%%%%%%%%%%%%%%%%%%%%%%%%%%%%%%%%%%%%%%%%%%%%%%%
    %%% Ajustes para a classe EXAM:
    \input{/app/static/utils/exam}

    %%%%%%%%%%%%%%%%%%%%%%%%%%%%%%%%%%%%%%%%%%%%%%%%%%%%%%%%%%%%%%%%%%%%%%%%%%%%%%%%
    %%% Ajustes para textos:
    \input{/app/static/utils/texto}

    %%%%%%%%%%%%%%%%%%%%%%%%%%%%%%%%%%%%%%%%%%%%%%%%%%%%%%%%%%%%%%%%%%%%%%%%%%%%%%%%
    %%% Ambientes, aliases, comandos e customizações em geral para serem
    %%% utilizadas nos documentos. Defina aqui qualquer customização específica
    %%% para seus documentos!
    \input{/app/static/utils/customizacoes}

    %%%%%%%%%%%%%%%%%%%%%%%%%%%%%%%%%%%%%%%%%%%%%%%%%%%%%%%%%%%%%%%%%%%%%%%%%%%%%%%%
    %%% Ajuste do layout, espaçamento de linhas e etc.:
    \geometry{a4paper, portrait, left=2.0cm, right=2.0cm, top=2.0cm, bottom=2.0cm}
    \singlespacing

    %%%%%%%%%%%%%%%%%%%%%%%%%%%%%%%%%%%%%%%%%%%%%%%%%%%%%%%%%%%%%%%%%%%%%%%%%%%%%%%%
    %%% Configurações de cabeçalho e rodapé (não use fancyhdr pois ocorre conflito):
    \pagestyle{headandfoot}
    \runningheadrule
    \firstpageheader{}{}{}
    \runningheader{Sem disciplina}{}{Avaliação}
    \firstpagefooter{}{}{Página \thepage\ de \numpages}
    \runningfooter{}{}{Página \thepage\ de \numpages}
    \runningfootrule

    %%%%%%%%%%%%%%%%%%%%%%%%%%%%%%%%%%%%%%%%%%%%%%%%%%%%%%%%%%%%%%%%%%%%%%%%%%%%%%%%
    %%% Configurações para as propriedades do PDF:
    \hypersetup{
    hidelinks,           % Comente para web, descomente para impressão
    colorlinks=true,      % True para web, False para impressão
    pdftitle=teste: Avaliação,
    pdfauthor=bren,
    pdfsubject=Sem disciplina,
    pdfkeywords=['Sem tag'],
    pdfinfo={
        CreationDate={}, % Ex.: D:AAAAMMDDHH24MISS
        ModDate={}       % Ex.: D:AAAAMMDDHH24MISS
    }
    }
    \input{/app/static/utils/pdfconfig}

    %%%%%%%%%%%%%%%%%%%%%%%%%%%%%%%%%%%%%%%%%%%%%%%%%%%%%%%%%%%%%%%%%%%%%%%%%%%%%%%%
    %%% Começa o documento
    \begin{document}

    %%%%%%%%%%%%%%%%%%%%%%%%%%%%%%%%%%%%%%%%%%%%%%%%%%%%%%%%%%%%%%%%%%%%%%%%%%%%%%%%
    %%% Folha de instruções padronizadas da UVV para a prova
    \pdfbookmark[1]{Instruções}{instrucoes}
    \begin{figure}[H]
        \begin{center}
            \includegraphics[scale=0.3]{{/app/static/imgs/logo-uvv.png}}
        \end{center}	
    \end{figure}

    \vspace{-1cm}

    \begin{table}[H]
        \centering
        \begin{tabular}{|llll|}
            \hline
            \multicolumn{2}{|l|}{\textbf{Disciplina:} Sem disciplina} & 
            \multicolumn{1}{l|}{\multirow{3}{*}{\textbf{\makecell{Nota:\\ \,\\ \,}}}} & 
            \multirow{3}{*}{\textbf{\makecell{Coordenador:\\ \,\\ \,}}} \\ 
            \cline{1-2}
            \multicolumn{2}{|l|}{\textbf{Professor:} bren \hspace{4.72cm} } & 
            \multicolumn{1}{l|}{} & \\ 
            \cline{1-2}
            \multicolumn{2}{|l|}{\textbf{Aluno:}} & 
            \multicolumn{1}{l|}{} & \\ 
            \hline
            \multicolumn{1}{|l|}{\textbf{Turma:} \hspace{6.3cm} } & 
            \multicolumn{1}{l|}{\textbf{Semestre:}} & 
            \multicolumn{2}{l|}{\textbf{Valor:} 10 pontos} \\ 
            \hline
            \multicolumn{1}{|l|}{\textbf{Data:}} & 
            \multicolumn{3}{l|}{\textbf{Avaliação:}} \\ 
            \hline
        \end{tabular}
    \end{table}

    \begin{center}
        \fbox{\setlength\fboxsep{0.3cm} \fbox{\parbox{15.4cm}{
            \begin{center}
                \textbf{INSTRUÇÕES PARA A REALIZAÇÃO DESTA AVALIAÇÃO:}
            \end{center}
            \begin{itemize}[leftmargin=*]
                \item Esta é a primeira avaliação da disciplina Sem disciplina, 
                    e engloba o conteúdo referente Sem tag.
                \item Esta avaliação contém 10 questões objetivas, \textbf{todas obrigatórias}.
                \item \textbf{Leia atentamente cada questão!} As questões objetivas terão uma,
                    e apenas uma única, resposta a ser marcada.
                \item \textbf{Desligue o celular} antes de começar. A avaliação é
                    \textbf{individual} e \textbf{sem consulta}.
                \item As questões podem ser respondidas, na prova, com lápis ou caneta. Ao final
                    da prova \textbf{{VOCÊ DEVERÁ PREENCHER O GABARITO COM CANETA
                    DE COR PRETA}} \textbf{\underline{OU AZUL}} para que seja feita a correção
                    automatizada através da leitura do padrão de respostas marcado no
                    gabarito.
                \item \textbf{CUIDADO ao preencher o gabarito} pois eles são \textbf{INDIVIDUAIS
                    e NOMEADOS} (cada estudante tem um gabarito próprio específico, com um
                    código de identificação único). \textbf{Questões rasuradas são anuladas}
                    pelo software de leitura do gabarito.
                \item Ao final da prova, \textbf{DEVOLVA A PROVA E O GABARITO} para o professor.
                \item Siga todas as normas de \textbf{Integridade Acadêmica} da disciplina.
                    Alunos que forem flagrados com qualquer espécie de ``cola'' ou trocando
                    informações com outros alunos terão suas avaliações recolhidas, as notas
                    zeradas, e a situação será encaminhada para a coordenação para a aplicação
                    das penalidades previstas pela universidade.
                \item Com 90 minutos de avaliação você terá tempo de até 9,min por questão,
                    em média.
                \item Boa avaliação!
            \end{itemize}
            \vspace{2cm}
        }}}
    \end{center}
    \newpage

    %%%%%%%%%%%%%%%%%%%%%%%%%%%%%%%%%%%%%%%%%%%%%%%%%%%%%%%%%%%%%%%%%%%%%%%%%%%%%%%%
    %%% Inicia o bloco de questões:
    \begin{questions}

    \newpage
    \question
Analise o diagrama relacional do banco de dados exibido na figura abaixo:
\begin{figure}[H]
  \centering
  \caption{Banco de dados de peças e fornecedores}
  \label{fig:pecas}
  %\fbox{
    \includegraphics[width=0.5\linewidth]{/app/static/imgs/defaul.png}
  %}\\
  %\footnotesize{Fonte: xxxx.}
\end{figure}
Podemos dizer que esse banco de dados ilustra:
\begin{choices}
\choice Um relacionamento com cardinalidade N:N entre as tabelas ``pecas' e ``fornecimentos', realizado através de relacionamentos não identificados com a tabela ``fornecedores'. 
\choice Um relacionamento com cardinalidade N:N entre as tabelas ``pecas' e ``fornecedores', realizado através de relacionamentos identificados através da tabela ``fornecimentos'. 
\choice Um relacionamento com cardinalidade N:N entre as tabelas ``pecas' e ``fornecedores', realizado através de relacionamentos não identificados através da tabela ``fornecimentos'. 
\choice Um relacionamento com cardinalidade 1:N entre as tabelas ``pecas' e ``fornecedores', indicando que uma peça pode ter sido fornecida por diversos fornecedores. 
\choice Um relacionamento com identificação N:N entre as tabelas ``pecas' e ``fornecedores', realizado através do relacionamento com cardinalidade identificada através da tabela ``fornecimentos'. 
\end{choices}

\question
Filtro da mediana é:
\begin{choices}
\choice Indicado para atenuar ruído com preservação de bordas (i.e., rápidas transições de nível em imagens).
\choice Indicado para detectar bordas em imagens.
\choice Indicado para detectar tonalidades específicas em uma imagem.
\choice Indicado para detectar formas específicas em imagens.
\choice Nenhuma das respostas acima.
\end{choices}

\question
Quando trabalhando com sistemas baseados em trocas de mensagens, temporizações (\textit{time-outs}) são utilizadas para:

\begin{choices}
\choice Temporariamente suspender a transmissão de mensagens.
\choice Limitar o tempo para obter um recurso.
\choice Arbitrar que uma mensagem transmitida foi perdida.
\choice Limitar o tamanho de uma mensagem transmitida.
\choice Limitar o número de retransmissões de uma mensagem.
\end{choices}

\question
Considere a seguinte tabela em uma base de dados relacional (chave primária sublinhada):\\
Tabela1 (CodAluno, CodDisciplina, AnoSemestre, NomeAluno, NomeDisciplina, CodNota, DescricaoNota)\\
Considere as seguintes dependências funcionais:
\begin{itemize}
    \item CodAluno \(\rightarrow\) NomeAluno
    \item CodDisciplina \(\rightarrow\) NomeDisciplina
    \item (CodAluno, CodDisciplina, AnoSemestre) \(\rightarrow\) CodNota
    \item (CodAluno, CodDisciplina, AnoSemestre) \(\rightarrow\) DescricaoNota
    \item CodNota \(\rightarrow\) DescricaoNota
\end{itemize}

Considerando as formas normais, qual das afirmativas abaixo se aplica:
\begin{choices}
\choice A tabela está na quarta forma normal.
\choice A tabela encontra-se na primeira forma normal, mas não na segunda forma normal.
\choice A tabela encontra-se na terceira forma normal, mas não na quarta forma normal.
\choice A tabela não está na primeira forma normal.
\choice A tabela encontra-se na segunda forma normal, mas não na terceira forma normal.
\end{choices}

\question Assinale a alternativa \textbf{incorreta}:
\begin{choices}
\choice O MariaDB é um SGBD
\choice Modelo de dados é a mesma coisa que o diagrama relacional
\choice Um SGBD não tem interface gráfica
\choice No modelo relacional, tudo o que o usuário percebe são tabelas
\choice O modelo relacional de dados foi criado por Edgar Frank Cood
\end{choices}

\question
(Adaptado de ENADE 2011) Considere o diagrama abaixo, que ilustra um banco de
dados que descreve os empregados que trabalham nos departamentos de uma empresa:
\begin{figure}[H]
  \centering
  \caption{Empregados e departamentos}
  \vspace{-0.2cm}
  \label{fig:empregados}
  %\fbox{
    \includegraphics[width=0.5\linewidth]{/app/static/imgs/defaul.png}
  %}\\
  %\footnotesize{Fonte: xxxx.}
\end{figure}
\vspace{-0.4cm}
Na linguagem SQL, o comando mais simples para recuperar os códigos dos
departamentos que contém empregados que ganham menos do que 2000 e mais do que
5000 reais é:

\begin{choices}
\choice \begin{verbatim}SELECT cod_departamento
FROM empregados
WHERE salario <= 2000 AND salario >= 5000;
\end{verbatim}

\choice \begin{verbatim}SELECT cod_departamento
FROM empregados
WHERE salario BETWEEN 2000 AND 5000;
\end{verbatim}

\choice \begin{verbatim}SELECT cod_departamento
FROM empregados
WHERE salario < 2000 AND salario > 5000;
\end{verbatim}

\choice \begin{verbatim}SELECT cod_departametno
FROM empregados
WHERE salario <= 2000 OR salario >= 5000;
\end{verbatim}

\choice \begin{verbatim}SELECT cod_departamento
FROM empregados
WHERE salario < 2000 OR salario > 5000;
\end{verbatim}

\end{choices}

\question
Professor Mac Sperto propôs o seguinte algoritmo de ordenação, chamado de Super Merge, similar ao \textit{merge sort}: divida o vetor em 4 partes do mesmo tamanho (ao invés de 2, como é feito no \textit{merge sort}). Ordene recursivamente cada uma das partes e depois intercale-as por um procedimento semelhante ao procedimento de intercalação do \textit{merge sort}. Qual das alternativas abaixo é verdadeira?
\begin{choices}
\choice Super Merge não está correto. Não é possível ordenar quebrando o vetor em 4 partes
\choice Super Merge está correto, mas consome tempo maior que \( O(\text{merge sort}) \)
\choice Nenhuma das afirmações acima está correta
\choice Super Merge está correto, mas consome tempo menor que \( O(\text{merge sort}) \)
\choice Super Merge está correto, mas consome tempo \( O(\text{merge sort}) \)
\end{choices}

\question
Na Álgebra Booleana:
\begin{choices}
\choice Há dez valores motivados pelos dez dedos do ser humano
\choice Os dígitos são alfanuméricos que podem ser representados por um byte
\choice Os dígitos são hexadecimais de 0 a 15
\choice Os dígitos são octais, de 0 a 7
\choice Os dígitos são binários 0 e 1
\end{choices}

\question
Seja a árvore binária abaixo a representação de um espaço de estados para um problema \( p \), em que o estado inicial é \( a \), e \( i \) e \( f \) são estados finais.

\begin{figure}[H]
    \centering
    \includegraphics[width=0.5\linewidth]{/app/static/imgs/defaul.png}
    \caption{Questão 59}
    \label{fig:enter-label}
\end{figure}

Um algoritmo de busca em largura-primeiro forneceria a seguinte sequência de estados como primeira alternativa a um caminho-solução para o problema \( p \):
\begin{choices}
\choice \( a \ b \ d \ h \ e \ i \)
\choice \( a \ b \ e \ i \)
\choice \( a \ b \ d \ e \ f \)
\choice \( a \ b \ c \ d \ e \ f \)
\choice \( a \ c \ f \)
\end{choices}

\question
Considere o projeto de um circuito digital que implementa uma função \( f \) com três variáveis de entrada e satisfazendo as seguintes propriedades:

\[
f(x, y, z) = 
\begin{cases} 
1 & \text{se } x \neq y \\ 
0 & \text{caso contrário} 
\end{cases}
\]

Qual das seguintes expressões representa corretamente a função \( f \)?
\begin{choices}
\choice \( xy + \overline{y}z + \overline{z} \)
\choice \( \overline{x}z + xy + yz \)
\choice \( \overline{x}y z + x y z \)
\choice \( x + \overline{y}z \)
\choice \( \overline{x}y + \overline{y} \)
\end{choices}



    %%%%%%%%%%%%%%%%%%%%%%%%%%%%%%%%%%%%%%%%%%%%%%%%%%%%%%%%%%%%%%%%%%%%%%%%%%%%%%%%
    %%% Finaliza o bloco de questões:
    \end{questions}
    \end{document}
    